\documentclass[12 pt, twoside]{amsart}

\usepackage{../current_style}

\begin{document}
\title{Almond Cookies}
\maketitle
\subtitle{20 minutes of prep and 20 minutes to bake.}

\noindent \desc{These are a family favorite, and it's no surprise. They don't require much, take little time to prepare, and taste really good. Even more, it's naturally dairy free because it doesn't even call for oil. 
Between the almond extract, lemon, and spices it can be easy to overpower the cookies with spices, but it's not like they will taste bad if you do ;). The powdered sugar can be substituted for a bit less than a cup of normal sugar.}

\recipe{
\begin{ingredients}
  \item $2$ $\frac{1}{4}$ cups almond flour
  \item $1$ $\frac{3}{4}$ cups powdered sugar*
  \item $\frac{1}{4}$ tsp baking powder
  \item (optional) 1 tsp cardamom
  \item (optional) 1 tsp nutmeg
  \item whites of two eggs
  \item zest of one lemon
  \item 1 tbsp lemon juice
  \item 1 tsp almond extract
  \item 2 tsp vanilla extract
  \item Decorations
\end{ingredients}
}{
\begin{directions}
  \item[1.] Preheat oven to 300 Fahrenheit. Mix all dry ingredients
  \item[2a.] Get the egg whites from two eggs and beat the whites until there are soft peaks, OR
  \item[2b.] Whisk the egg whites from two eggs until they have soft peaks. 
  \item[3.] Add the extracts and lemon juice to the egg and gently whip. Add to the dry ingredients and incorporate.
  \item[4.] Roll into small balls, roll in powdered sugar, and flatten. If you want, top with something decorative like chocolate or almonds, then bake for 20 minutes or until the bottoms are browning.
\note{powdered sugar can be substituted for a bit less than a cup of normal sugar. I would definitely put less in though if you are going to roll them in sugar before baking.}
\end{directions}
}


\end{document}
